\clearpage
\hypertarget{native-270-pure17}{}\label{native-270-pure17}
\Article{17}
Zur weiteren Verdeutlichung greifen wir das Beispiel aus Artikel~13 wieder auf, in dem wir
\[
\begin{aligned}
A'=A'''&=2+4R+R^3+2R^4=2R-2R^2-R^3\\
A''&=2+R+4R^2+2R^3=-R+2R^2-2R^4
\end{aligned}
\]
gefunden haben. Nehmen wir für $R$ den Wert $\cos72^0+i\sin72^0$, so wird
\[
\begin{aligned}
A'=A'''&=2\cos72^0-3\cos144^0+i(2\sin72^0-\sin144^0)\\
A''&=-3\cos72^0+2\cos144^0+i(\sin72^0+2\sin144^0)
\end{aligned}
\]
Die Winkel $b',b''$ werden daher durch die Gleichungen
\[
\begin{aligned}
1)\qquad\sin b'&=\frac{2\sin72^0-\sin144^0}{\surd11}\\
2)\qquad\cos b'&=\frac{2\cos72^0-3\cos144^0}{\surd11}\\
3)\qquad\operatorname{tang}b'&=\frac{2\sin72^0-\sin144^0}{2\cos72^0-3\cos144^0}\\
4)\qquad\sin b''&=\frac{\sin72^0+2\sin144^0}{\surd11}\\
5)\qquad\cos b''&=\frac{-3\cos72^0+2\cos144^0}{\surd11}\\
6)\qquad\operatorname{tang}b''&=\frac{\sin72^0+2\sin144^0}{-3\cos72^0+2\cos144^0}
\end{aligned}
\]
bestimmt. Jede der Gleichungen $1,2,3$ genügt, um den Winkel $b'$ zu bestimmen, sofern der Quadrant bekannt ist, in dem er zu nehmen ist; dieser muss anhand der Vorzeichen der Größen $2\sin72^0-\sin144^0$, $2\cos72^0-3\cos144^0$ bestimmt werden. Dasselbe gilt für den Winkel $b''$. In unserem Fall ist $b'$ zwischen $0$ und $90^0$, $b''$ zwischen $90^0$ und $180^0$ zu nehmen. Multiplizieren wir Zähler und Nenner von Gleichung $3$ mit $-3\cos72^0+2\cos144^0$, so geht sie in
\SourceLeaf{271}
\[
\operatorname{tang}b'=\frac2{31}\{-\sin72^0+13\sin144^0\}
\]
über. Ebenso ergibt sich aus Gleichung $6$, wenn Zähler und Nenner mit $2\cos72^0-3\cos144^0$ multipliziert werden,
\[
\operatorname{tang}b''=\frac2{31}\{-13\sin72^0-\sin144^0\}
\]
In Zahlen erhalten wir daraus
\[
\begin{aligned}
\operatorname{tang}b'&=+0{,}4316226944,\quad\log\operatorname{tang}b'=9{,}6351042715\\
&\hspace{1em}b'=23^020'46''04603\\
\operatorname{tang}b''&=-0{,}8355819332,\quad\log\operatorname{tang}b''=9{,}9219890411\mathrm n\\
&\hspace{1em}b''=140^07'6''52441
\end{aligned}
\]
woraus sich
\[
5f'=186^048'38''61647,\qquad f'=37^021'43''723294
\]
ergibt.

.Wir haben also
\[
\begin{aligned}
(2,1)&=-\frac15+\frac{\surd11}5\{2\cos37^021'43''723294+2\cos51^022'41''400558\}\\
(2,2)&=-\frac15+\frac{\surd11}5\{2\cos325^021'43''723294+2\cos267^022'41''400558\}\\
(2,4)&=-\frac15+\frac{\surd11}5\{2\cos253^021'43''723294+2\cos123^022'41''400558\}\\
(2,8)&=-\frac15+\frac{\surd11}5\{2\cos181^021'43''723294+2\cos339^022'41''400558\}\\
(2,5)&=-\frac15+\frac{\surd11}5\{2\cos109^021'43''723294+2\cos195^022'41''400558\}
\end{aligned}
\]
woraus wir
\[
\begin{array}{rclcl}
(2,1)&=&+1{,}6825070652&=&2\cos\frac{360^0}{11}\\[3pt]
(2,2)&=&+0{,}8308299&=&2\cos\frac{720^0}{11}\\[3pt]
(2,4)&=&&=&2\cos\frac{1440^0}{11}\\[3pt]
(2,8)&=&&=&2\cos\frac{2880^0}{11}\\[3pt]
(2,5)&=&&=&2\cos\frac{1800^0}{11}
\end{array}
\]
finden.

\Needspace{13\baselineskip}
\Article{18}
\begingroup
\setlength{\abovedisplayskip}{8pt plus 2pt minus 2pt}
\setlength{\belowdisplayskip}{8pt plus 2pt minus 2pt}
Ein weiteres Beispiel liefert uns die Gleichung $x^{17}-1=0$, die wir in den \emph{Disquis. Arithm.} schon nach einer anderen Methode behandelt haben. Setzen wir also $n=17$, $\beta=8$, $g=3$; dann entsprechen den Zahlen die folgenden Indizes:
\SourceLeaf{272}
\[
\begin{array}{r*{16}{@{\;}r}}
\text{Zahlen}&1\pdot&2\pdot&3\pdot&4\pdot&5\pdot&6\pdot&7\pdot&8\pdot&9\pdot&10\pdot&11\pdot&12\pdot&13\pdot&14\pdot&15\pdot&16\\
\text{Indizes}&0\pdot&14\pdot&1\pdot&12\pdot&5\pdot&15\pdot&11\pdot&10\pdot&2\pdot&3\pdot&7\pdot&13\pdot&4\pdot&9\pdot&6\pdot&8
\end{array}
\]
Daraus finden wir
\[
\begin{aligned}
A'=A''''''&=2R+2R^2+3R^4+4R^5+2R^6+2R^7\\
A''=A'''''&=2+3R+R^3+R^4+3R^5+4R^6+R^7\\
A'''=A''''&=3+3R+2R^2+3R^3+R^5+2R^6+R^7
\end{aligned}
\]
oder, da in diesem Fall $R^4+1=0$ gilt,
\[
\begin{aligned}
A'=A''''''&=-3-2R-2R^3\\
A''=A'''''&=1-4R^2\\
A'''=A''''&=3+2R+2R^3
\end{aligned}
\]
Setzen wir also $R=\cos45^0+i\sin45^0$, so wird
\[
\begin{aligned}
A'=A''''''&=-3-2i\surd2,\\
A''=A'''''&=1-4i,\\
A'''=A''''&=3+2i\surd2
\end{aligned}
\]

Somit finden wir $b',b'',b'''$ aus den Gleichungen
\[
\begin{aligned}
\sin b'&=-\surd\tfrac8{17}&\sin b''&=-\surd\tfrac{16}{17}&\sin b'''&=+\surd\tfrac8{17}\\
\cos b'&=-\surd\tfrac9{17}&\cos b''&=+\surd\tfrac1{17}&\cos b'''&=+\surd\tfrac9{17}\\
\operatorname{tang}b'&=+\surd\tfrac89&\operatorname{tang}b''&=-4&\operatorname{tang}b'''&=+\surd\tfrac89
\end{aligned}
\]
woraus wir
\[
b'=223^018'49'',\quad b''=284^02'10,\quad b'''=43^018'49''=b'-180^0
\]
\[
4f'=550^039'48'',\qquad f'=137^039'57''
\]
\begingroup\postdisplaypenalty=10000
\[
\!\begin{aligned}
(2,1)&=\mkern-2mu-\frac18+\frac{\surd17}8\{2\cos137^039'57''+2\cos52^01'5''+2\cos265^038'52''+1\}\\
(2,4)&=\mkern-2mu-\frac18+\frac{\surd17}8\{2\cos92^039'57''+2\cos322^01'5''+2\cos130^038'52''-1\}\\
(2,9)&=\mkern-2mu-\frac18+\frac{\surd17}8\{2\cos47^039'57''+2\cos232^01'5''+2\cos355^038'52''+1\}\\
(2,10)&=\mkern-2mu-\frac18+\frac{\surd17}8\{2\cos2^039'57''+2\cos142^01'5''+2\cos220^038'52''-1\}\\
(2,13)&=\mkern-2mu-\frac18+\frac{\surd17}8\{2\cos317^039'57''+2\cos52^01'5''+2\cos85^038'52''+1\}\\
(2,5)&=\mkern-2mu-\frac18+\frac{\surd17}8\{2\cos272^039'57''+2\cos322^01'5''+2\cos310^038'52''-1\}\\
(2,15)&=\mkern-2mu-\frac18+\frac{\surd17}8\{2\cos227^039'57''+2\cos232^01'5''+2\cos175^038'52''+1\}\\
(2,11)&=\mkern-2mu-\frac18+\frac{\surd17}8\{2\cos182^039'57''+2\cos142^01'5''+2\cos40^038'52''-1\}
\end{aligned}
\]
herleiten.
\endgroup

\SourceLeaf{273}
\[
\begin{array}{rclcl}
\frac12(2,1)&=&+0{,}092268&=&\cos\frac4{17}360^0\\[3pt]
\frac12(2,3)&=&&=&\cos\frac5{17}360^0\\[3pt]
\frac12(2,9)&=&&=&\cos\frac2{17}360^0\\[3pt]
\frac12(2,10)&=&&=&\cos\frac6{17}360^0\\[3pt]
\frac12(2,13)&=&+0{,}93247&=&\cos\frac1{17}360^0\\[3pt]
\frac12(2,5)&=&&=&\cos\frac3{17}360^0\\[3pt]
\frac12(2,15)&=&&=&\cos\frac8{17}360^0\\[3pt]
\frac12(2,11)&=&&=&\cos\frac7{17}360^0
\end{array}
\]
\endgroup
\begin{center}*\qquad\qquad*\\*\end{center}

\Needspace{21\baselineskip}
Von diesen allgemeineren Untersuchungen über die Funktionen $[r,R]$, welche die in Artikel~360 der \emph{Disquiss. Ar.} begonnene zweite Theorie der reinen Gleichungen weiter verdeutlichen und erweitern, gehen wir zu einer genaueren Betrachtung gewisser besonderer Fälle über, nämlich solcher, in denen für $\beta$ bestimmte Werte genommen werden. Daraus werden mehrere ebenso fruchtbare wie elegante Untersuchungen hervorgehen; einige davon waren zwar schon in den \emph{Disq. Ar.} (Artikel\ .\ .\ .), jedoch nach einer anderen Methode, behandelt worden, andere sind dagegen als völlig neu anzusehen. Den bemerkenswerten Zusammenhang dieser Untersuchungen mit der höheren Arithmetik, die dadurch sehr große und bisher unerwartete Fortschritte macht, wollen wir in einer anderen, bald zu veröffentlichenden Abhandlung darlegen. — Ferner setzen wir in der gesamten folgenden Untersuchung voraus, dass für $r$ eine \emph{eigentliche} Wurzel der Gleichung $x^n-1=0$ und für $R$ eine eigentliche Wurzel der Gleichung $R^\beta-1=0$ genommen wird.

\Article{19}
Wir beginnen mit dem Wert $\beta=2$, wobei also für $R$ der Wert $-1$ zu nehmen ist. Unsere Funktion $[r,R]$ wird daher
\[
=r-r^g+r^{g^2}-r^{g^3}\dotsfour-r^{g^{n-2}}
\]
und es gilt
\[
[r,R]=-[r^g,R]=+[r^{g^2},R]=-[r^{g^3},R]\quad\text{usw.}
\]
Allgemein gilt, wenn $\lambda$ eine beliebige, nicht durch $n$ teilbare ganze Zahl bezeichnet,
\SourceLeaf{274}
\[
\begin{aligned}
\relax[r^\lambda,R]&=+[r,R]\quad\text{wenn }\lambda\text{ quadratischer Rest von }n\text{ ist},\\
\relax[r^\lambda,R]&=-[r,R]\quad\text{wenn }\lambda\text{ quadratischer Nichtrest von }n\text{ ist}.
\end{aligned}
\]
Bezeichnen wir ferner die quadratischen Reste von $n$, die unter $1,2,3\dots n-1$ enthalten sind, unbestimmt mit $a$ und die Nichtreste von $n$ innerhalb derselben Grenzen mit $b$, so sind offenbar die Zahlen
\[
1,g^2,g^4\dotsfour g^{n-3}
\]
ohne Rücksicht auf ihre Reihenfolge zum Modul $n$ den Zahlen $a$ kongruent, und ebenso die Zahlen
\[
g,g^3,g^5\dotsfour g^{n-2}
\]
den Zahlen $b$, sodass $[r,R]=\Sigma r^a-\Sigma r^b$ wird.

\Needspace{15\baselineskip}
Setzen wir also $\frac{360^0}n=\omega$ und $r=\cos k\omega+i\sin k\omega$, so wird $[r,R]=\Sigma\cos ak\omega-\Sigma\cos bk\omega+i\Sigma\sin ak\omega-i\Sigma\sin bk\omega$. Nach Artikel~14 ist das Quadrat der Funktion $[r,R]$ gleich $=+n$ oder $=-n$, je nachdem, ob $n$ von der Form $4z+1$ oder $4z-1$ ist; somit gilt im ersten Fall $[r,R]=\pm\surd n$, im zweiten $[r.R]=\pm i\surd n$. Das der Wurzelgröße vorangestellte Vorzeichen bleibt jedoch mehrdeutig. Daraus ergeben sich die folgenden Summen:

\noindent I.\quad Ist $n$ von der Form $4z+1$, so gilt
\[
\begin{aligned}
\Sigma\cos ak\omega-\Sigma\cos bk\omega&=\pm\surd n\\
\Sigma\sin ak\omega-\Sigma\sin bk\omega&=0
\end{aligned}
\]
\noindent II.\quad Ist $n$ von der Form $4z-1$, so gilt
\[
\begin{aligned}
\Sigma\cos ak\omega-\Sigma\cos bk\omega&=0\\
\Sigma\sin ak\omega-\Sigma\sin bk\omega&=\pm\surd n
\end{aligned}
\]

Da ferner die Gesamtheit der Zahlen $a,b$ offenbar mit $1,2,3\dots n-1$ übereinstimmt, gilt $\Sigma r^a+\Sigma r^b=r+r^2+r^3+\text{usw.}+r^{n-1}=-1$, und damit $\Sigma\cos ak\omega+\Sigma\cos bk\omega=-1$, $\Sigma\sin ak\omega+\Sigma\sin bk\omega=0$. Aus diesen und den vorigen Summen folgen die nachstehenden:

\noindent I.\quad Für den ersten Fall
\[
\begin{aligned}
\Sigma\cos ak\omega&=-\tfrac12\pm\tfrac12\surd n\\
\Sigma\cos bk\omega&=-\tfrac12\mp\tfrac12\surd n\\
\Sigma\sin ak\omega&=\Sigma\sin bk\omega=0
\end{aligned}
\]

\SourceLeaf{275}
\noindent II.\quad Für den zweiten Fall
\[
\begin{aligned}
\Sigma\cos ak\omega&=\Sigma\cos bk\omega=-\tfrac12\\
\Sigma\sin ak\omega&=\pm\tfrac12\surd n\\
\Sigma\sin bk\omega&=\mp\tfrac12\surd n
\end{aligned}
\]
\Needspace{9\baselineskip}
Diese Summen wurden bereits in den \emph{Disquiss. Arr.}, Artikel~356, nach einer nicht sehr verschiedenen Methode hergeleitet; keine der beiden Methoden kann jedoch die Mehrdeutigkeit des der Wurzelgröße voranzustellenden Vorzeichens beseitigen. Diesen Mangel haben wir aber vor Kurzem in einer besonderen Abhandlung behoben, in der bewiesen wurde, dass für den Wert $k=1$ in allen angeführten Formeln die oberen Vorzeichen zu nehmen sind.
\BookRule{.3\textwidth}
